% =====================================================================
%  Chapter 19 — Disk-Resolved Imaging and the Close-Approach Window
% =====================================================================

\chapter{Disk-Resolved Imaging and the Close-Approach Window}
\label{ch:disk-imaging}

There are two ways to see an asteroid as a disk rather than as a point. One
is to be near it. The other is to have an aperture large enough, or a
baseline long enough, that its angular size is resolved from far away. The
first is the method that has actually produced images; the second is the
method this book is proposing. This chapter is about the first, and about
what it costs.

\section{What a resolved disk gives}

A resolved image of a rotating body is worth more than a lightcurve, and the
reason is geometric. A lightcurve gives the body's brightness as a single
number at each time. A resolved image gives the brightness \emph{distribution}
across the body, and from that distribution the shape follows directly.

Two effects are immediately visible in a disk image and invisible in a
lightcurve.

The \emph{terminator}, the boundary between the day and night sides, shows
the orientation of the body directly: the line across the disk from the
brightest point to the darkest is the projection of the spin axis, and no
model is needed to find it.

The \emph{limb profile}, the shape of the boundary between the disk and the
background, shows the shape of the body's cross-section at that epoch. A
concave body shows a concave limb.

This is why the space missions that have imaged an asteroid from close range
have produced the field's best shape models, and why the number of such
images is the field's most quoted scarcity: six.

\section{Adaptive optics on the ground}

Ground-based adaptive optics\idx{Adaptive optics} removes the atmospheric turbulence that blurs an
image, and it extends the resolution of a large telescope to the diffraction
limit of its aperture. For a bright near-Earth object at a close approach,
an eight-metre telescope with adaptive optics\idx{Adaptive optics} can resolve a body of a few
hundred metres across at a distance of a tenth of an astronomical unit.

The limitation is not the resolution but the window, and the window has two
parts.

The first is the proximity, and it is unavoidable: the object's apparent
size must exceed the resolution of the instrument, so the distance is
bounded by the object's size.

The second is the phase angle. The phase angle\idx{Phase angle} is the angle at the object
between the direction to the Sun and the direction to the observer. At a
close approach, the phase angle\idx{Phase angle} changes rapidly, and the illumination
geometry determines which part of the body is visible. An observation made
at a high phase angle\idx{Phase angle}, when the terminator\idx{Terminator} runs across the disk and most of
the visible face is in shadow, yields a shape model with a large systematic
uncertainty; an observation made at a low phase angle, near full
illumination, yields a clean one.

\begin{deeper}{Why the wrong phase angle wastes the encounter}
The illumination geometry at a close approach is not under the observer's
control in the usual way. It is fixed by the orbit. A conjunction geometry,
in which the object passes between the Earth and the Sun, is ideal for
imaging and poor for radar; an opposition geometry, in which the object is
behind the Earth, is the reverse. The lesson is not that a close approach
should be sought, but that a close approach should be \emph{anticipated}, so
that the geometry is known in advance and the instruments can be scheduled
for it. That is an argument for an automated alert path, not merely for a
telescope.
\end{deeper}

\section{What the phase angle does to the inversion}

The sensitivity of the shape inversion to the phase angle is the reason the
literature is cautious about the value of a single close-range image. The
brightness distribution across a disk depends on the illumination geometry
as well as on the shape, and the two are not separable from a single
observation at a single phase angle.

The practical consequence is that a disk image is normally combined with a
lightcurve, and the pair is inverted together, with the lightcurve
constraining the rotation and the disk constraining the shape. Where the
phase angle is favourable, the disk dominates the shape constraint; where it
is not, the lightcurve dominates, and the answer reverts to a convex model.

This is a further argument for a measurement that does not depend on the
illumination geometry at all. An interferometric measurement of a rotating
body measures the silhouette, and the silhouette is a geometric boundary
rather than a brightness distribution. It is insensitive to the phase angle
in a way that a photometric inversion is not.

\section{The scarcity, and what it costs}

\begin{table}[!htbp]
\centering
\small
\caption{The resolved-shape record, and what each entry cost.}
\label{tab:resolved-record}
\begin{tabularx}{\textwidth}{@{}P{1.43in}P{1.13in}Y@{}}
\toprule
\rowcolor{panel}
\textbf{Method} & \textbf{Models} & \textbf{What it costs} \\
\midrule
Spacecraft fly-by & six images & a multi-billion-dollar mission per target \\
\rowcolor{panel}
Ground adaptive optics & a handful & a chance encounter at the right phase \\
Lightcurve inversion & about 25,000 & a convex hull \\
\rowcolor{panel}
Radar & a few dozen & a radar-crossing target, and proximity \\
\bottomrule
\end{tabularx}
\end{table}

The costs in \tabref{tab:resolved-record} are what define the opportunity.
A fly-by image costs a mission. A ground adaptive-optics image costs a
chance. A radar image costs a radar-crossing target and proximity. A
lightcurve costs a telescope time, and returns a convex hull.

There is no row in that table for a measurement that costs a telescope time,
returns a resolved shape, and requires no proximity. That is the gap the
book's proposal is aimed at, and the argument of Part V is that the
lunar surface is the only place where that row can be filled.

\begin{keytakeaways}
\item A resolved disk gives the orientation and the limb profile directly,
both of which a lightcurve cannot supply.
\item The binding constraint on ground imaging is the phase angle, not the
resolution, and the right phase has to be anticipated rather than hoped for.
\item Six close-range images exist, some hundreds of radar models, and about
25,000 convex lightcurve inversions.
\item The gap is a measurement that costs telescope time, returns a resolved
shape, and needs no proximity.
\end{keytakeaways}
